% !TEX root = ../../main.tex
% File: part2/chapters_part2/chap2_5.tex
% Cập nhật 2026: thực hành hậu huấn luyện với TRL

\section{Huấn luyện Hậu kỳ với TRL: SFT, DPO và GRPO \newtag}
\label{sec:posttraining_trl}

Phần 1 đã trình bày lý thuyết của tinh chỉnh có giám sát (chương~\ref{chap:finetuning_alignment}), căn chỉnh theo sở thích (mục~\ref{sec:alignment_methods}) và học tăng cường cho suy luận (chương~\ref{chap:reasoning}). Trong thực tế, cả ba giai đoạn này đều có thể được thực hiện với thư viện \textbf{TRL} (Transformer Reinforcement Learning) của Hugging Face -- được xây dựng trên \texttt{transformers}, \texttt{accelerate} và \texttt{peft}, và tích hợp vLLM để sinh văn bản nhanh trong quá trình RL.

\begin{tcolorbox}[title={Quy trình hậu huấn luyện điển hình}, colback=blue!5!white, colframe=blue!60!black, fonttitle=\bfseries]
\textbf{Mô hình nền} $\xrightarrow{\text{SFT}}$ \textbf{mô hình tuân theo chỉ dẫn} $\xrightarrow{\text{DPO/ORPO/KTO}}$ \textbf{mô hình được căn chỉnh} $\xrightarrow{\text{GRPO + phần thưởng kiểm chứng được}}$ \textbf{mô hình suy luận tốt hơn trên toán/mã}.

Không phải dự án nào cũng cần đủ ba bước: nhiều bài toán doanh nghiệp chỉ cần SFT (thường với LoRA/QLoRA); GRPO chỉ đáng giá khi bạn có \emph{hàm thưởng tự động đáng tin cậy}.
\end{tcolorbox}

\subsection{Định dạng dữ liệu}
\label{ssec:trl_data_formats}
TRL chấp nhận cả định dạng văn bản thường lẫn định dạng hội thoại (danh sách các tin nhắn có \texttt{role} và \texttt{content}); với định dạng hội thoại, chat template của mô hình được áp dụng tự động.

\begin{example}{Các định dạng dữ liệu cho SFT, DPO và GRPO}{ex:trl_formats}
\begin{minted}{python}
# SFT -- hội thoại (tính loss trên phần trả lời)
{"prompt": [{"role": "user", "content": "Thủ đô của Pháp?"}],
 "completion": [{"role": "assistant", "content": "Paris."}]}

# DPO -- dữ liệu sở thích: cùng prompt, một câu được chọn, một câu bị loại
{"prompt": [{"role": "user", "content": "Bầu trời màu gì?"}],
 "chosen": [{"role": "assistant", "content": "Màu xanh."}],
 "rejected": [{"role": "assistant", "content": "Màu xanh lá."}]}

# GRPO -- chỉ cần prompt; các cột khác (ví dụ đáp án) được
# truyền vào hàm thưởng dưới dạng tham số cùng tên
{"prompt": [{"role": "user", "content": "12 * 7 = ?"}],
 "ground_truth": "84"}
\end{minted}
\end{example}

\subsection{Bước 1 -- SFT với LoRA và QLoRA}
\label{ssec:trl_sft}

\begin{example}{SFT với LoRA bằng SFTTrainer}{ex:trl_sft_lora}
\begin{minted}{python}
# pip install trl peft datasets
from datasets import load_dataset
from peft import LoraConfig
from trl import SFTConfig, SFTTrainer

dataset = load_dataset("trl-lib/Capybara", split="train")

peft_config = LoraConfig(
    r=16,
    lora_alpha=32,                # hệ số tỷ lệ = alpha / r = 2
    lora_dropout=0.05,
    target_modules="all-linear",  # gắn LoRA vào mọi lớp tuyến tính
    task_type="CAUSAL_LM",
)

training_args = SFTConfig(
    output_dir="qwen3-0.6b-sft-lora",
    learning_rate=1e-4,           # LoRA dùng LR lớn hơn full fine-tuning
    num_train_epochs=1,
    per_device_train_batch_size=4,
    gradient_accumulation_steps=4,
    max_length=2048,
    packing=True,                 # ghép nhiều mẫu ngắn vào một chuỗi
    assistant_only_loss=True,     # chỉ tính loss trên lời của assistant
)

trainer = SFTTrainer(
    model="Qwen/Qwen3-0.6B",
    args=training_args,
    train_dataset=dataset,
    peft_config=peft_config,
)
trainer.train()
trainer.save_model()
\end{minted}
\end{example}

\paragraph{QLoRA: tinh chỉnh mô hình lớn trên một GPU.} Chỉ cần thêm cấu hình lượng tử hóa 4-bit NF4 (mục~\ref{ssec:reparam}); mô hình gốc được lưu ở 4 bit, adapter LoRA vẫn huấn luyện ở BF16.

\begin{example}{QLoRA với SFTTrainer}{ex:trl_qlora}
\begin{minted}{python}
import torch
from transformers import BitsAndBytesConfig

bnb_config = BitsAndBytesConfig(
    load_in_4bit=True,
    bnb_4bit_quant_type="nf4",          # NormalFloat 4-bit
    bnb_4bit_use_double_quant=True,     # lượng tử kép
    bnb_4bit_compute_dtype=torch.bfloat16,
)

trainer = SFTTrainer(
    model="Qwen/Qwen3-8B",
    args=training_args,
    train_dataset=dataset,
    peft_config=peft_config,
    quantization_config=bnb_config,
)
\end{minted}
\end{example}

\subsection{Bước 2 -- Căn chỉnh theo sở thích với DPO}
\label{ssec:trl_dpo}
\begin{example}{DPO với DPOTrainer}{ex:trl_dpo}
\begin{minted}{python}
from datasets import load_dataset
from peft import LoraConfig
from trl import DPOConfig, DPOTrainer

dataset = load_dataset("trl-lib/ultrafeedback_binarized", split="train")

training_args = DPOConfig(
    output_dir="qwen3-0.6b-dpo",
    beta=0.1,               # càng lớn càng bám sát mô hình tham chiếu
    learning_rate=1e-5,     # (với LoRA); full fine-tuning thường ~5e-7
    per_device_train_batch_size=4,
    gradient_accumulation_steps=8,
    max_length=2048,
    # loss_type=["sigmoid"] là DPO gốc; TRL còn hỗ trợ "ipo",
    # "sigmoid_norm" (SimPO), "robust", ... để thử nghiệm
)

trainer = DPOTrainer(
    model="qwen3-0.6b-sft-lora",   # xuất phát từ mô hình đã SFT
    args=training_args,
    train_dataset=dataset,
    peft_config=LoraConfig(target_modules="all-linear"),
)
trainer.train()
\end{minted}
\end{example}
Khi không truyền \texttt{ref\_model}, TRL dùng chính mô hình ở trạng thái trước khi huấn luyện làm mô hình tham chiếu $\pi_{\text{ref}}$ (với LoRA, chỉ cần tắt adapter là có mô hình tham chiếu, không tốn thêm bộ nhớ). Các chỉ số cần theo dõi: \texttt{rewards/margins} (khoảng cách phần thưởng ngầm giữa câu được chọn và bị loại -- nên tăng dần) và \texttt{rewards/accuracies}. TRL cũng cung cấp \texttt{KTOTrainer} cho dữ liệu nhãn nhị phân và hỗ trợ ORPO (mục~\ref{sec:alignment_methods}).

\subsection{Bước 3 -- Học tăng cường với GRPO và phần thưởng kiểm chứng được}
\label{ssec:trl_grpo}
Ví dụ dưới đây huấn luyện một mô hình nhỏ giải toán GSM8K. Hàm thưởng gồm hai phần: \textbf{thưởng định dạng} (suy nghĩ trong thẻ \texttt{<think>}, đáp án trong \texttt{\textbackslash boxed\{\}}) và \textbf{thưởng chính xác} (so khớp đáp số) -- đúng tinh thần của DeepSeek-R1-Zero (mục~\ref{ssec:r1_zero}).

\begin{example}{GRPO trên GSM8K với hàm thưởng dựa trên luật}{ex:trl_grpo}
\begin{minted}{python}
import re
from datasets import load_dataset
from trl import GRPOConfig, GRPOTrainer

SYSTEM = ("Suy nghĩ từng bước trong <think>...</think>, "
          "sau đó ghi đáp số cuối cùng trong \\boxed{}.")

def to_prompt(example):
    # Đáp án của GSM8K nằm sau chuỗi "####"
    answer = example["answer"].split("####")[-1].strip().replace(",", "")
    return {
        "prompt": [{"role": "system", "content": SYSTEM},
                   {"role": "user", "content": example["question"]}],
        "ground_truth": answer,
    }

dataset = load_dataset("openai/gsm8k", "main", split="train").map(to_prompt)

def _text(completion):
    # completion ở dạng hội thoại: [{"role": "assistant", "content": ...}]
    return completion[0]["content"]

def format_reward(completions, **kwargs):
    pattern = re.compile(r"<think>.*?</think>.*\\boxed\{.+?\}", re.DOTALL)
    return [1.0 if pattern.search(_text(c)) else 0.0 for c in completions]

def accuracy_reward(completions, ground_truth, **kwargs):
    rewards = []
    for c, gt in zip(completions, ground_truth):
        found = re.findall(r"\\boxed\{([^}]*)\}", _text(c))
        pred = found[-1].strip().replace(",", "") if found else None
        rewards.append(1.0 if pred == gt else 0.0)
    return rewards

training_args = GRPOConfig(
    output_dir="qwen3-0.6b-grpo-gsm8k",
    learning_rate=1e-6,
    num_generations=8,             # G: số câu trả lời mỗi prompt
    max_completion_length=1024,
    per_device_train_batch_size=8,
    beta=0.0,                      # không dùng KL -> không cần mô hình tham chiếu
    loss_type="dapo",              # chuẩn hóa theo token (mục DAPO)
    epsilon_high=0.28,             # Clip-Higher
    use_vllm=True,                 # sinh văn bản nhanh bằng vLLM
    vllm_mode="colocate",          # vLLM chạy chung GPU với trainer
    logging_steps=10,
)

trainer = GRPOTrainer(
    model="Qwen/Qwen3-0.6B",
    reward_funcs=[format_reward, accuracy_reward],
    reward_weights=[0.2, 1.0],
    args=training_args,
    train_dataset=dataset,
)
trainer.train()
\end{minted}
\end{example}
Chạy trên nhiều GPU bằng \texttt{accelerate launch train\_grpo.py}. Các tham số liên hệ trực tiếp với lý thuyết:
\begin{itemize}
	\item \texttt{num\_generations} là kích thước nhóm $G$; nhóm quá nhỏ làm ước lượng advantage nhiễu.
	\item \texttt{loss\_type}: \texttt{"grpo"} (chuẩn hóa theo từng câu trả lời như bài báo gốc), \texttt{"dapo"} (chuẩn hóa theo token), \texttt{"dr\_grpo"} (bỏ thiên lệch độ dài) -- xem mục~\ref{ssec:grpo_variants}.
	\item \texttt{beta}: hệ số KL; nhiều công thức hiện đại đặt bằng 0 để tiết kiệm bộ nhớ.
	\item \texttt{scale\_rewards}: có chia cho độ lệch chuẩn của nhóm hay không (Dr.~GRPO khuyến nghị không).
\end{itemize}

\begin{tcolorbox}[title={Kinh nghiệm thực chiến khi chạy GRPO}, colback=red!5!white, colframe=red!75!black, fonttitle=\bfseries]
\begin{itemize}
	\item \textbf{Kiểm tra hàm thưởng trước khi huấn luyện:} Chạy mô hình gốc trên vài chục prompt, in ra câu trả lời và phần thưởng. Hàm thưởng lỗi (ví dụ regex quá chặt) khiến mọi phần thưởng bằng 0 và mô hình không học gì.
	\item \textbf{Chọn độ khó phù hợp:} Nếu mô hình gốc giải đúng gần 0\% hoặc gần 100\% số bài, hầu hết các nhóm có advantage bằng 0. Lọc dữ liệu về vùng mô hình giải đúng một phần.
	\item \textbf{Theo dõi:} phần thưởng trung bình (nên tăng), độ dài câu trả lời, tỷ lệ bị cắt ngưỡng (clip ratio), entropy (giảm quá nhanh là dấu hiệu sụp đổ khám phá).
	\item \textbf{Cảnh giác reward hacking:} Đọc mẫu đầu ra định kỳ -- mô hình rất giỏi tìm kẽ hở của hàm thưởng (ví dụ in nhiều \texttt{\textbackslash boxed\{\}} để “trúng” đáp án).
	\item \textbf{Đánh giá độc lập:} Đo trên tập kiểm tra riêng bằng framework đánh giá chuẩn (mục~\ref{sec:lm_eval_practice}), không dùng phần thưởng huấn luyện làm thước đo.
\end{itemize}
\end{tcolorbox}
